\documentclass[11pt]{article}
\usepackage[T1]{fontenc}
\usepackage{textcomp}
\usepackage[margin=0.75in]{geometry}
\usepackage{amsmath,amssymb,amsthm}
\usepackage{array,booktabs}
\usepackage{hyperref}
\setlength{\parindent}{0pt}
\newtheorem{theorem}{Theorem}
\theoremstyle{definition}
\newtheorem{definition}{Definition}
\title{Flow Proof Artifact: prop-29-equal-parallelograms-same-base}
\author{Generated by \texttt{flow doc proof}}
\date{Flow Proof Book}
\begin{document}
\maketitle
In equal parallelograms which have the same base the sides opposite the base are in the same straight lines.\\[0.5em]
\textbf{Source.} euclid — Elements, Book VI, Proposition 29\\[0.5em]
\bigskip
\noindent{\large \textbf{Derived fact 1.} \textit{In equal parallelograms which have the same base the sides opposite the base are in the same straight lines} \hfill the Euclidean plane \cdot Euclid Book VI \cdot Proposition 29: equal parallelograms on the same base are in the same parallels}\par
\smallskip\noindent\textit{Coordinate: the Euclidean plane \cdot Euclid Book VI \cdot Proposition 29: equal parallelograms on the same base are in the same parallels}\par
\smallskip\noindent\textit{Source: euclid — Elements, Book VI, Proposition 29}\par
\smallskip\noindent\textit{Built on: proposition 39: equal triangles on same base and same side are in same parallels, for Book I of the Elements in on the Euclidean plane}\par
\medskip
\noindent\textbf{Goal.} In equal parallelograms which have the same base the sides opposite the base are in the same straight lines\par
\noindent$\displaystyle \text{equal pars same base same parallels}$\par
\medskip
\noindent
\begin{tabular}{@{} >{\raggedright\arraybackslash}p{0.06\textwidth} >{\raggedright\arraybackslash}p{0.47\textwidth} >{\raggedleft\arraybackslash}p{0.41\textwidth} @{}}
\textbf{\#} & \textbf{Proof} & \textbf{Mathematics} \\
\midrule
\textbf{1.} & We prove that proposition 29: equal parallelograms on the same base are in the same parallels for Euclid Book VI the Euclidean plane. & \\[0.45em]
\textbf{2.} & Let ABCD = parallelogram\_1. & \\[0.45em]
\textbf{3.} & Let ABCE = parallelogram\_2. & \\[0.45em]
\textbf{4.} & Let base\_AB = common\_base. & \\[0.45em]
\textbf{5.} & From step 2, step 3, and step 4, we can deduce that in same parallels. & $\displaystyle \text{in same parallels}$ \\[0.45em]
\textbf{6.} & We invoke the derived fact governing Book I of the Elements in the Euclidean plane: proposition 39: equal triangles on same base and same side are in same parallels, for Book I of the Elements in on the Euclidean plane. & \\[0.45em]
\textbf{7.} & From step 6, this implies equal pars same base same parallels. Hence proven. & $\displaystyle \text{equal pars same base same parallels}$ \\[0.45em]
\bottomrule
\end{tabular}
\label{thm:Geometry-Euclid Book VI-proposition_29_equal_parallelograms_on_the_same_base_are_in_the_same_parallels}
\par\medskip\hrule
\end{document}
